\note{sa page 5}

façon équivalente, il existe un $v \in \mathcal{E}^{(2n-2i)}$, tel que le morphisme
induit $v \colon M^{i} \to M^{2n-i}$ soit un isomorphisme.

On peut transformer ceci de la façon suivante :

\[ (4.4)_{i} \qquad \begin{cases} \exists\, v_{i} \in \mathcal{E}^{(2n-2i)},\ w_{i} \in \mathcal{E}^{-(2n-2i)} \ \text{tels que l'on ait} \\[2pt] (w_{i} v_{i} - 1)\, \pi_{i} = 0 , \qquad (v_{i} w_{i} - 1)\, \pi_{2n-i} = 0 \end{cases} \]

Or, sans plus supposer que $\mathcal{E}^{*}$ se réalise à l'aide d'un $M^{*}$,
on voit aussitôt que $(4.4)_{i}$ implique $(3.1)_{i}$ et $(3.2)_{i}$. En effet,
quitte à remplacer $v$, $w$ par $v \pi_{i} = \pi_{2n-i} v$,
$w \pi_{2n-i} = \pi_{i} w$, on peut supposer $v$, $w$ dans
$\mathcal{E}^{2n-i}_{i}$, $\mathcal{E}^{i}_{2n-i}$. Si alors
\struck{$u \in \mathcal{E}^{(0)}$} $u \in \mathcal{E}^{(0)}$ est tel que
$\mathcal{E}^{2n-i}_{i} u = 0$, donc $v u = 0$, on aura
$w v u \pi_{i} = 0$, i.e.\ $w v \pi_{i} u = 0$, donc $\pi_{i} u = 0$ (car
$w v \pi_{i} = \pi_{i}$). De même pour $(3.2)_{i}$, en utilisant la dernière
relation $(4.4)_{i}$.

\marginal{\struck{et 5.} v. p. 5 bis}

6. Nous allons étudier de façon plus détaillée le cas où les conditions (4.4)
sont satisfaites. Pour ceci, \struck{considérons} $\mathcal{E}^{*}$ \add{étant
une sous-algèbre graduée de $\operatorname{End}(M^{*})$} \ill{} $M^{*}$,
\struck{et supposons} qu'on ait un

\[ (6.0) \qquad L \in \mathcal{E}^{(2)} \]

qui ait la propriété que

